%% kapa: exam=1 round=19 year=2008 period=1 form=A subject=civil_law subject_ko=민법 seq=10 number=10
%% source: Q-Net 감정평가사 기출문제 (cst00309) artlSeq=1001962 file=extracted/1008/1001962/2008년 제19회 감정평가사 1차시험 기출문제/1교시.hwp
%% conversion: kordoc 4.24.51 -> tools/gen1.py (auto); review=proofread-sonnet; flags=none
\begin{question}{10}
부재자의 재산관리인에 관한 설명으로 옳은 것은? (다툼이 있는 경우에는 판례에 의함)
\begin{choices}
  \item 재산관리인에 의한 부재자 소유의 부동산 매매행위에 대한 법원의 허가결정은 과거의 매매를 추인하는 방법으로는 할 수 없다.
  \item 재산관리인이 무이자부 금전소비대차를 이자부 소비대차로 하는 행위는 법원의 허가를 얻어야 한다.
  \item 법원에 의해 재산관리인의 선임이 있은 후에, 부재자가 사망한 것으로 확실해진 경우에는 관리인의 권한은 당연히 소멸한다.
  \item 부재자의 재산관리인의 대리권의 범위는 친권자의 대리권의 범위와 같다.
  \item 법원이 선임한 재산관리인은 법정대리인이지만, 부재자 자신이 선임한 재산관리인은 임의대리인이다.
\end{choices}
\end{question}
\answer{5}
